%% ma: L12-B08-0030
%% lop: 12
%% chuong: 2
%% bai: 8
%% dang: 12.8.21
%% loai: TN4
%% muc: NB
%% dap_an: "A"
%% nguon: "LaTeX giáo viên gửi 21/10/2026 (M05, Vectơ và hệ trục tọa độ trong không gian, 107 câu), Phần I Câu 3"
%% kien_thuc: "Bài 8 (tọa độ vectơ khi biết tọa độ hai điểm)"
%% trang_thai: cho-duyet
%% ly_do: "Dạng mới đề xuất 12.8.21"
\begin{ex}
Trong không gian $Oxyz$, cho hai điểm $A(1;-2;3)$ và $B(3;0;-1)$. Tọa độ của vectơ $\overrightarrow{AB}$ là
\choice{\True $(2;2;-4)$}{$(4;-2;2)$}{$(2;2;4)$}{$(2;-2;-4)$}
\loigiai{$\overrightarrow{AB}=(3-1;0+2;-1-3)=(2;2;-4)$.}
\end{ex}
